
Quality filtering is a cornerstone of modern LLM data curation. Perplexity-based selection, which retains documents scoring well against reference language models trained on high-quality corpora, has emerged as a preferred approach in large-scale training pipelines. However, this filtering mechanism may introduce systematic bias: benchmark content---educational text, Q\&A forums, multiple-choice questions---exhibits the structured, low-perplexity characteristics that such filters preferentially select.

This work investigates the coupling between data curation and benchmark contamination. While prior work has developed contamination detection methods (n-gram overlap, membership inference) and studied curation effects on downstream performance, no prior work has connected these domains to measure how curation decisions affect contamination rates.

The central insight is that data attribution methods can bridge contamination detection and curation analysis. By combining contamination detection (identifying which examples overlap with benchmarks) with attribution methods (measuring how much each example contributes to benchmark performance), one can quantify the Contamination Contribution Ratio (CCR): the fraction of benchmark-specific attribution mass originating from contaminated examples.

\subsubsection{Contributions}

This work presents:

\begin{enumerate}
\item \textbf{Contamination Contribution Ratio (CCR)}: A metric combining n-gram contamination detection with TRAK attribution to quantify benchmark-specific influence from contaminated training examples. Synthetic injection experiments validate CCR as calibrated ($R^2$ = 0.9998).
\end{enumerate}

\begin{enumerate}
\item \textbf{Evidence of curation-contamination coupling}: In methodology validation experiments with simulated contamination, perplexity filtering increases CCR by 0.1594 relative to random sampling (p < 0.0001).
\end{enumerate}

\begin{enumerate}
\item \textbf{Causal validation via removal intervention}: Removing high-CCR examples causes 1.$97\times$ greater accuracy degradation than random removal (95\% CI: [1.53, 2.34]), indicating contaminated examples are causally necessary for benchmark performance.
\end{enumerate}

\begin{enumerate}
\item \textbf{Influence Fragility Ratio (IFR)}: Contaminated examples exhibit 4.$1\times$ higher IFR than non-contaminated examples (p < $10^{-162})$, suggesting structural irreplaceability. However, the IFR-redundancy correlation ($\rho$ = -0.11) was weaker than the hypothesized threshold ($\rho$ < -0.5).
\end{enumerate}
